\subsubsection{MLP Variants}
\label{app:exp:mlp-variant-definitions}

Across the standalone fact-storage-capacity experiments (\Cref{sec:fact-storage-capacity,app:exp:section4:capacity}) and the Transformer-block capacity experiments (\Cref{appendix:lm-scaling}), we compare the following MLP variants at matched hidden width $m$ on the same synthetic fact sets:

\begin{description}
    \item[\textbf{GD (gradient-descent trained).}] Our default GD baseline is a bilinear gated-identity MLP, i.e.\ the bilinear specialization of the gated MLP family in \Cref{eq:gated-mlp-class}.
    Unless stated otherwise, we train these models for
    $10{,}000$ epochs using Adam with initial learning rate $10^{-3}$ and a
    cosine-annealing schedule down to $10^{-6}$.

    \item[\textbf{NTK.}] Our NTK baseline is the degree-$1$ Hermite weight construction of~\citet{nichani2024understandingfactualrecalltransformers}, which uses a gated ReLU MLP architecture.
    A self-contained description is given in \Cref{app:theory:hebbian_mlps:ntk_construction}.

    \item[\textbf{Our construction.}] Our closed-form construction is a Hebbian MLP with
    sketched-$K_{2}$ kernel, i.e. the bilinear random-feature construction described in
    \Cref{sec:sketched-k2,alg:sketched_k2_hebbian_whitening}.
    We study three variants:
    \begin{itemize}
        \item \textbf{Unwhitened.} This is the raw sketched-$K_{2}$ Hebbian construction,
        with the bilinear random-feature map used directly in the final Hebbian readout.
        \item \textbf{Whitened.} This uses the same bilinear feature map, but applies the
        full readout whitening procedure from
        \Cref{app:theory:margin_bounds:kernel_whitening} with ridge parameter $10^{-6}$.
        \item \textbf{Data-dependent.} This keeps the same bilinear architecture, but performs least squares solves to obtain the random feature vectors within the Hebbian kernel, as described in
        \Cref{app:theory:margin_bounds:data_dependent_construction}, then applies whitening with ridge parameter $10^{-6}$.
    \end{itemize}
\end{description}

\subsubsection{Transformer Setups}
\label{app:exp:transformer-setups}
\label{app:ssfr-training-setup}
We describe details about the modified 1-layer, 1-head GPT-style Transformer architecture used across all of our SSFR experiments unless otherwise specified.

We freeze the input and output embeddings, and when an experiment includes a fact-storing MLP, that MLP uses the same embeddings as the Transformer.
Positional encodings are disabled throughout.
We use RMSNorm before the attention layer, and use unit-RMSNorm (which projects to the unit sphere) before the MLP and before the final Transformer output.
We train with AdamW using learning rate $2\times 10^{-4}$, batch size $1{,}280$, and $4{,}000$ iterations.

Our experiments include two different training protocols:
\begin{itemize}
    \item In the \emph{pretrained attention} setting, we first train attention alone on a dummy SSFR task using the identity map, so that it learns to query the MLP when faced with distractor junk tokens in its context.
    After training, we freeze the attention layer, then train a frozen fact-storing MLP on a fresh random-permutation fact set using the same embeddings, insert it into the Transformer, and evaluate the combined model on SSFR with that fact set.
    \item In the \emph{inserted MLP} setting, we first construct or train a frozen fact-storing MLP on a random-permutation fact set, then insert it into the Transformer block. We then train the attention layer around that frozen MLP on SSFR using the same fact set.
\end{itemize}
